\section{Examples of fields}
\label{section-examples}

\noindent
To get started, let us begin by providing several examples of fields. The
reader should recall that if $R$ is a ring and $I \subset R$ an
ideal, then $R/I$ is a field precisely when $I$ is a maximal ideal.

\begin{example}[Rational numbers]
\label{example-rational-numbers}
The rational numbers form a field. It is called the
{\it field of rational numbers} and denoted $\mathbf{Q}$.
\end{example}

\begin{example}[Prime fields]
\label{example-prime-field}
If $p$ is a prime number, then $\mathbf{Z}/(p)$ is a field, denoted
$\mathbf{F}_p$. Indeed, $(p)$ is a
maximal ideal in $\mathbf{Z}$. Thus, fields may be finite: $\mathbf{F}_p$
contains $p$ elements.
\end{example}

\begin{example}
\label{example-quotient-polymial-ring}
In a principal ideal domain, an ideal generated by an irreducible element
is maximal. Now, if $k$ is a field, then the polynomial ring $k[x]$ is a
principal ideal domain. It follows that if $P \in k[x]$ is an irreducible
polynomial (that is, a nonconstant polynomial
that does not admit a factorization into terms of smaller degrees), then
$k[x]/(P)$ is a field. It contains a copy of $k$ in a natural way.
This is a very general way of constructing fields. For instance, the
complex numbers $\mathbf{C}$
can be constructed as $\mathbf{R}[x]/(x^2 + 1)$.
\end{example}

\begin{example}[Quotient fields]
\label{example-quotient-field}
Recall that, given a domain $A$, there is an imbedding $A \to F$ into a
field $F$ constructed from $A$ in exactly the same manner that
$\mathbf{Q}$ is constructed from $\mathbf{Z}$. Formally the elements
of $F$ are (equivalence classes of) fractions $a/b$,
$a, b \in A$, $b \not = 0$. As usual $a/b = a'/b'$ if and only if $ab' = ba'$.
The field $F$ is called the {\it quotient field}, or {\it field of fractions},
or {\it fraction field} of $A$.
The quotient field has the following universal property: given an
injective ring map $\varphi : A \to K$ to a field $K$, there is a unique
map $\psi : F \to K$ making
$$
\xymatrix{
F \ar[r]_\psi & K \\
A \ar[u] \ar[ru]_\varphi
}
$$
commute. Indeed, it is clear how to define such a map: we set
$\psi(a/b) = \varphi(a)\varphi(b)^{-1}$ where injectivity of $\varphi$
assures that $\varphi(b) \not = 0$ if $ b \not = 0$.
\end{example}

\begin{example}[Field of rational functions]
\label{example-field-of-rational-functions}
If $k$ is a field, then we can consider the field $k(x)$ of rational
functions over $k$. This is the quotient field of the polynomial ring
$k[x]$. In other words, it is the set of quotients $F/G$ for
$F, G \in k[x]$, $G \not = 0$ with the obvious equivalence relation.
\end{example}

\begin{example}
\label{example-field-of-meromorphic-functions}
Let $X$ be a Riemann surface. Let $\mathbf{C}(X)$ denote the
set of meromorphic functions on $X$. Then $\mathbf{C}(X)$ is a ring under
multiplication and addition of functions. It turns out that in fact
$\mathbf{C}(X)$ is a field. Namely, if a nonzero function $f(z)$ is
meromorphic, so is $1/f(z)$. For example, let $S^2$ be the Riemann
sphere; then we know from complex analysis that the ring of meromorphic
functions $\mathbf{C}(S^2)$ is the field of rational functions $\mathbf{C}(z)$.
\end{example}



